\documentclass[10pt,a4paper]{article}
\usepackage[utf8]{inputenc}
\usepackage{geometry}
\geometry{top=1.5cm, bottom=1.5cm, left=1.8cm, right=1.8cm}
\usepackage{titlesec}
\titleformat{\section}{\small\bfseries\scshape}{}{0pt}{}[\titlerule]
\titlespacing*{\section}{0pt}{6pt}{4pt}
\usepackage{enumitem}
\setlist{noitemsep,topsep=2pt,parsep=0pt,partopsep=0pt}
\pagestyle{empty}

\begin{document}

\begin{center}
    {\large \textbf{COURSE COMPLETION PROJECT BRIEF: 1-PAGE RESEARCH CHARTER}}\\[2pt]
    \textbf{Course:} AI Engineering \quad | \quad \textbf{Supervisor:} PhD Academic Lead \quad | \quad \textbf{Date:} November 2026\\[4pt]
    \textbf{Student Felix ID:} 12 \quad | \quad \textbf{Degree Milestone:} Bachelor of Science Thesis (Computer Science, 5th Semester)
\end{center}

\vspace{-4pt}

\section{Working Titles}
\begin{enumerate}
    \item \textbf{Primary:} Quantization Versus Retrieval Precision: Evaluating Faithfulness and Abstention in Local Retrieval-Augmented Generation Over Computer Science Textbooks Under a 12 GB VRAM Constraint
    \item \textbf{Alternative A:} Resource-Constrained Educational RAG: How GGUF Weight Quantization and Hybrid Structural Retrieval Interact to Govern Hallucination and Abstention on Consumer GPUs
    \item \textbf{Alternative B:} Beyond Parametric Memorization: Contamination-Controlled Benchmarking of Quantized Open-Weight LLMs and Hybrid Retrieval Over Openly Licensed Programming Curricula
\end{enumerate}

\section{Project Idea \& Motivation}
Deploying Retrieval-Augmented Generation (RAG) assistants for university programming courses on consumer hardware (12 GB VRAM) requires aggressive weight quantization (4-bit to 8-bit GGUF) and compact open-weight models (3B--14B parameters). While general RAG chunking is well studied, two critical methodological flaws undermine existing local educational evaluations: (1) failure to control for pre-training parametric memorization of public textbooks, and (2) uncalibrated evaluation of how low-bit quantization interacts with retrieval noise to degrade claim-level faithfulness and unanswerable query abstention. This project constructs a reproducible, network-isolated (\texttt{--network none}) local RAG evaluation harness and releases a human-verified 300-item QA benchmark (255 answerable with gold evidence spans, 45 unanswerable abstention probes) over the CC BY 4.0 OpenStax textbook \textit{Introduction to Python Programming} to quantify the joint sensitivity of quantization bit-width and hybrid structural retrieval.

\section{Bounded Domain \& Scientific Delta (SOTA Alignment)}
\begin{itemize}
    \item \textbf{Target Domain Scope:} \textit{In-Scope:} OpenStax Python textbook (CC BY 4.0); 300-item verified QA benchmark; Study 1 retrieval ablation (fixed-size dense vs.\ AST/Markdown structure-aware dense vs.\ structure-aware hybrid BM25+dense with cross-encoder reranking) on an 8B Q4\_K\_M backbone; Study 2 quantization matrix (\{3B, 8B\} $\times$ \{Q4\_K\_M, Q5\_K\_M, Q8\_0\} + 14B Q4\_K\_M) vs.\ closed-book no-retrieval baseline on an RTX 4070 (12 GB VRAM, 8--12 GPU-hours total). \textit{Out-of-Scope:} Weight fine-tuning, multimodal OCR, $>12$ GB VRAM models (14B Q8\_0), and live classroom trials.
    \item \textbf{Scientific Baseline (SOTA):} Grounded against foundational RAG and IR benchmarks in WoS/Scopus/OpenAlex: Lewis et al.\ (2020), BEIR (Thakur et al., 2021), SQuAD 2.0 abstention (Rajpurkar et al., 2018), RAGAS (Es et al., 2024), ARES (Saad-Falcon et al., 2024), and post-training quantization literature (Frantar et al., 2023; Dettmers et al., 2023).
    \item \textbf{Expected Scientific Delta:} Isolates net retrieval gain ($\Delta \text{Correctness} = \text{Score}_{\text{RAG}} - \text{Score}_{\text{ClosedBook}}$) and demonstrates whether hybrid reranked retrieval compensates for 4-bit quantization loss on claim-level faithfulness and abstention F1, verified via NLI entailment (\texttt{DeBERTa-v3-MNLI}), human-calibrated LLM judging (Cohen's $\kappa \ge 0.70$), and paired bootstrap 95\% confidence intervals.
\end{itemize}

\section{Developmental Delta (Continuous Progress Rule)}
\begin{itemize}
    \item \textbf{Current Student Baseline:} Strong self-reported software baseline in Python, Git, Docker, Linux/Bash, SQL, Scikit-Learn, and PyTorch (9/10 self-assessment retained as documented baseline risk \texttt{R1}), but zero prior formal training in IR evaluation metrics, psychometric judge calibration, contamination controls, or LLMOps observability.
    \item \textbf{Specific Skill/Architecture Gaps to Bridge:} Master four new engineering pillars: (1) rank-aware IR metrics ($\text{Recall@}k, \text{MRR}, \text{nDCG}$) and claim-level NLI verification; (2) human annotation reliability (Cohen's $\kappa$ on 100 items) and paired bootstrap / Wilcoxon hypothesis testing with Holm-Bonferroni correction; (3) closed-book contamination control and abstention F1 benchmarking; and (4) reproducible LLMOps (\texttt{MLflow} experiment tracking, \texttt{Langfuse} tracing, pinned Ollama SHA-256 digests, and VRAM/latency profiling).
\end{itemize}

\section{Tooling Architecture \& Automated Pipelines}
\begin{itemize}
    \item \textbf{Agent Governance:} Google Antigravity IDE (\texttt{01\_ideation\_ceo}) + NotebookLM MCP for WoS/Scopus/OpenAlex synthesis and Zotero Better BibTeX citation management.
    \item \textbf{Compute \& Sandboxing:} Split execution across Linux laptop CPU (ingestion, embeddings, reranking, IR metrics, NLI verification) and University Lab Nvidia RTX 4070 12 GB VRAM (network-isolated Docker + Ollama batch generation, 8--12 GPU-hours across Weeks 4--10, with Google Colab Free T4 shard fallback), tracked in \texttt{MLflow} and \texttt{Langfuse}.
    \item \textbf{Artifact Archiving:} GitHub repository (\texttt{12W6PIE/Class.pro}) releases bridged to Zenodo (CERN Data Centre) minting a permanent DataCite DOI for the 300-item CC-BY QA benchmark, evaluation code, and raw generation logs.
\end{itemize}

\section{Course Deliverables Mapping}
\begin{itemize}
    \item \textbf{Track 1(a) AI Tech Utilization Report:} Comprehensive engineering manual detailing Docker \texttt{--network none} GPU sandboxing, GGUF VRAM budgeting, \texttt{MLflow}/\texttt{Langfuse} telemetry, NLI + LLM judge calibration, and Antigravity/MCP workflow orchestration.
    \item \textbf{Track 1(b) AI Tech Draft Manuscript:} Primary scientific manuscript reporting the two-study empirical evaluation of quantization $\times$ retrieval interactions on faithfulness, net RAG gain over closed-book memory, and unanswerable abstention.
    \item \textbf{Track 2(a) Domain Deliverables Report:} Dataset documentation, schema specification, and Zenodo DataCite DOI release notes for the 300-item OpenStax Python QA benchmark and reproducible evaluation harness.
    \item \textbf{Track 2(b) Domain Draft Manuscript:} Unified with Track 1(b) under the AI-Domain Exception Rule, incorporating educational CS error-taxonomy analysis (syntax vs.\ conceptual vs.\ code-tracing queries).
    \item \textit{AI Domain Exception Applied:} YES --- because the core research domain is AI Technology (resource-constrained RAG evaluation and quantization methodology), Track 1(a) and Track 1(b) fulfill the primary course mandate, supported by the Zenodo-archived benchmark artifact in Track 2(a).
\end{itemize}

\section{Statement of Publishing Level Assessment}
\textbf{Assessment:} Target **Tier 3 (International Conference Proceedings, e.g., OCETA Connect 2027 or IEEE/ACM applied AI \& CS education venues)**, with stretch potential for **Tier 2 (JCR Q3/Q4 indexed journal)** if the quantization $\times$ abstention interaction yields statistically significant cross-family findings ($\kappa \ge 0.70$, Holm-Bonferroni $p < 0.05$) and a reusable public benchmark. If empirical effect sizes between quantization tiers prove statistically indistinguishable on the 300-item sample, the rigorous negative/null methodology and Zenodo-archived dataset are guaranteed routing to **Tier 4 (Internal Semi-Scientific Course Publishing Service Database)**.

\end{document}
